\documentclass[11pt,a4paper]{article}

\usepackage[margin=2.8cm]{geometry}
\usepackage{amsmath,amssymb,amsthm,mathtools}
\usepackage{booktabs}
\usepackage{microtype}
\usepackage[hidelinks]{hyperref}
\usepackage[capitalize,noabbrev]{cleveref}

\newtheorem{theorem}{Theorem}
\newtheorem{lemma}[theorem]{Lemma}
\newtheorem{proposition}[theorem]{Proposition}
\newtheorem{corollary}[theorem]{Corollary}
\newtheorem{claim}[theorem]{Claim}
\theoremstyle{definition}
\newtheorem{remark}[theorem]{Remark}
\newtheorem{conjecture}[theorem]{Conjecture}
\newtheorem{definition}[theorem]{Definition}

\DeclareMathOperator{\crN}{cr}
\DeclareMathOperator{\chl}{\chi_{\ell}}
\DeclareMathOperator{\chDP}{\chi_{\mathrm{DP}}}
\DeclareMathOperator{\bw}{bw}
\newcommand{\NN}{\mathbb{N}}
\newcommand{\RR}{\mathbb{R}}
\newcommand{\eps}{\varepsilon}

\title{Two structural observations in the neighbourhood of\\
       Albertson's conjecture}
\author{(draft) Marc Lelarge\thanks{\texttt{marc.lelarge@gmail.com}.}}
\date{\today}

\begin{document}

\begin{center}
{\Large\bfseries WITHDRAWN (2026-05-17)}\\[1em]
\fbox{\parbox{0.92\textwidth}{\itshape
Observation~1 of this combined note is \textbf{false}. At
$t = 5$, a planar graph of Voigt~(1993) has $\chi_\ell(G) = 5$
(every planar graph is $5$-choosable by Thomassen~1994; Voigt's
example is not $4$-choosable) and $\crN(G) = 0 < 1 = \crN(K_5)$,
which directly falsifies the claim
``$\chi_\ell(G) \ge t$ and $t \le 18 \Rightarrow \crN(G) \ge
\crN(K_t)$''. Independently, the proof in this note appeals to
the Albertson--Cranston--Fox / Bar\'at--T\'oth / Ackerman chain
under the (incorrect) belief that the chain uses
$\chi$ only through Dirac's $\delta \ge t - 1$ bound;
Ackerman~(\texttt{arXiv:1509.01932}, \S 3.1) in fact uses the
minimum edge-count function $f_r(n)$ for $r$-critical graphs,
and the list-critical edge floor (Krivelevich~1997) is strictly
weaker, so the lift is unjustified even apart from the
$t = 5$ counterexample.\\[0.6em]
Observation~2 (the bisection-width Crossing Lemma for regular
spectral expanders) is salvageable, but the version here contains
several fixable but real errors: an odd-$n$ floor in Alon's
spectral bisection bound, a wrong rounding bookkeeping step in
\Cref{lem:PST}, a numerical illustration below the Bungener--Kaufmann
density threshold, and an overstated Ore-scope claim in the
limitations section. The corrected stand-alone spectral note is at
\texttt{D16\_expander\_crossing\_paper/expander\_crossing\_v2.tex}
(see \texttt{D17\_submission\_packets/paper\_D16.md} for the
current submission packet).\\[0.6em]
This combined-paper file is preserved for historical record only.
The companion list-coloring source
\texttt{D15\_list\_albertson\_paper/list\_albertson\_le\_18.tex}
is similarly withdrawn.
}}
\end{center}
\vspace{2em}

\maketitle

\begin{abstract}
We record two observations arising from an analysis of the
structural slacks in the chain of partial results towards
Albertson's conjecture. \textbf{First}, the
Albertson--Cranston--Fox~/ Bar\'at--T\'oth~/ Ackerman chain
establishing the conjecture unconditionally at chromatic number up
to $t = 18$ depends on its $\chi$-hypothesis only through the Dirac
minimum-degree bound, which lifts verbatim to list-coloring via the
list versions of Dirac's and Brooks' theorems (Borodin~1977;
Erd\H{o}s--Rubin--Taylor~1979). Assembling these observations gives
list-Albertson at $t \le 18$, strictly stronger than ordinary
Albertson at the same range. \textbf{Second}, the
Pach--Spencer--T\'oth bisection-width Crossing Lemma, combined with
Alon's spectral bisection lower bound, yields an explicit-$\theta$
Crossing Lemma for $d_0$-regular spectral $\theta$-expanders,
\[
   \crN(G) \;\ge\; \frac{(1 - \theta)^{2}\,d_0^{2}\,n^{2}}{1280}
                 \;-\; \frac{d_0^{2}\,n}{16},
\]
with an Albertson-type corollary on regular spectral-expander
critical graphs. We are not aware of an earlier Crossing-Lemma
inequality with explicit spectral dependence. Neither observation
closes the Cranston residual at $t \in \{25, 26\}$; both are
recorded as side-observations in the structural neighbourhood of
the conjecture.
\end{abstract}

\section{Introduction}\label{sec:intro}

Albertson's conjecture~\cite{Albertson} is the natural common
strengthening of the Four Colour Theorem and Guy's formula for
$\crN(K_n)$:

\begin{conjecture}[Albertson]\label{conj:albertson}
Every graph $G$ with $\chi(G) \ge t$ satisfies
$\crN(G) \ge \crN(K_t)$.
\end{conjecture}

The conjecture is currently known unconditionally for $t \le 24$
through a chain of partial results that all proceed by analysing a
minimum counterexample (MCE). Each step combines two ingredients:

\begin{enumerate}
  \item an \emph{edge floor} on critical graphs --
        $|E(G)| \ge \tfrac{k-1}{2}|V(G)|$ via Dirac's theorem on
        critical graphs~\cite{Dirac1952};
  \item a \emph{Crossing Lemma} input bounding $\crN(G)$ from below
        in terms of $|E(G)|^{3}/|V(G)|^{2}$ once $|E(G)| / |V(G)|$
        exceeds an explicit threshold.
\end{enumerate}

Together these inputs force a vertex bound $|V(G)| \le c t$ on any
hypothetical MCE; the small remaining range is then closed either
hand-by-hand or by a topological argument on graphs with few
crossings per edge. The chain proceeds in four steps:
Albertson--Cranston--Fox~\cite{ACF} (2010, $t \le 12$ at
$|V| \le 4t$), Bar\'at--T\'oth~\cite{BT} (2010, $t \le 16$ at
$|V| \le 3.57t$), Ackerman~\cite{Ackerman} (2019, $t \le 18$ at
$|V| \le 3.03t$ via $|E| \le 6n - 12$ for graphs with at most four
crossings per edge), and Cranston~\cite{Cranston} (2025, $t \le 24$
via a new structural lemma due to Fox--Pach--Suk~\cite{FPS}).

\paragraph{Two structural observations.}
This note records two observations on the structural neighbourhood
of this chain. Each is a "we observe that an existing chain of
theorems lifts or packages to one degree of generality higher"
statement; neither closes any open case of \Cref{conj:albertson}.

\paragraph{The list-coloring variant.}
A \emph{list assignment} on $G$ is a function $L : V(G) \to 2^{\NN}$;
a proper $L$-colouring is a proper colouring $c$ with $c(v) \in L(v)$
for every $v$. The graph $G$ is $L$-colourable if such a colouring
exists, and \emph{$k$-list-colourable} (or \emph{$k$-choosable}) if
it is $L$-colourable for every $L$ with $|L(v)| \ge k$ for all $v$.
The \emph{list chromatic number} $\chl(G)$ is the smallest such $k$.
The notion is due to Vizing~\cite{Vizing} and to Erd\H{o}s, Rubin
and Taylor~\cite{ERT}, and satisfies $\chl(G) \ge \chi(G)$ for every
graph $G$. The list version of Albertson's conjecture is
correspondingly stronger:

\begin{conjecture}[List-Albertson]\label{conj:list-albertson}
Every graph $G$ with $\chl(G) \ge t$ satisfies
$\crN(G) \ge \crN(K_t)$.
\end{conjecture}

Despite the conceptual closeness of the two conjectures, the
list-coloring variant appears to have received little attention in
the literature; for small $t \le 4$ it reduces to elementary
observations (any $G$ with $\chl(G) \ge 3$ is non-planar by the
list version of the Four Colour Theorem for planar graphs --
in fact every planar graph is $5$-choosable, by
Thomassen~\cite{Thomassen}). We are not aware of a result placing
list-Albertson at a threshold higher than $t = 4$.

\paragraph{The spectral Crossing-Lemma variant.}
On the topological side, the standard Crossing Lemma of Ajtai,
Chv\'atal, Newborn, Szemer\'edi~\cite{ACNS} and
Leighton~\cite{Leighton83} bounds $\crN(G)$ from below by
$c\,|E(G)|^{3}/|V(G)|^{2}$ once $|E(G)| \ge 4|V(G)|$; successive
improvements due to Pach--T\'oth~\cite{PachToth97},
Pach--Radoi\v{c}i\'c--Tardos--T\'oth~\cite{PRTT},
Ackerman~\cite{Ackerman}, and most recently Bungener and
Kaufmann~\cite{BK24} (at $c = 1/27.48$) have driven the constant
down. Parallel to this density-iteration line, Leighton~\cite{Leighton84}
introduced a bisection-width variant; Pach--Spencer--T\'oth~\cite{PachSpencerToth}
(via the dual estimate of Pach--Shahrokhi--Szegedy~\cite{PachShahrokhiSzegedy}
and S\'ykora--Vrt'o~\cite{SykoraVrto}) established
\begin{equation}\label{eq:PST-intro}
   \crN(G) \;\ge\; \frac{\bw(G)^{2}}{40} \;-\; \frac{\sum_v \deg(v)^{2}}{16},
\end{equation}
which we work with throughout in its self-contained $1/80$ form
(see \Cref{sec:PST}). Combining~\eqref{eq:PST-intro} with Alon's
spectral bisection bound for regular spectral expanders gives an
explicit dependence of the Crossing-Lemma constant on the spectral
gap parameter $\theta$.

\subsection{Statements}\label{ssec:statements}

The two main theorems of this note are the following.

\begin{theorem}[Observation 1: list-Albertson at $t \le 18$]\label{thm:main-list}
Let $G$ be a graph and let $t \le 18$ be an integer. If
$\chl(G) \ge t$, then $\crN(G) \ge \crN(K_t)$.
\end{theorem}

\begin{theorem}[Observation 2: bisection-width Crossing Lemma for
                regular spectral expanders]\label{thm:main-spec}
Let $G$ be a $d_0$-regular graph on $n$ vertices with second
adjacency eigenvalue $|\lambda_2(G)| \le \theta\,d_0$ for some
$\theta \in [0, 1)$. Then
\begin{equation}\label{eq:Tone-prime}
   \crN(G) \;\ge\;
     \frac{(1 - \theta)^{2}\,d_0^{2}\,n^{2}}{1280}
     \;-\; \frac{d_0^{2}\,n}{16}.
\end{equation}
The constant $1/1280$ can be replaced by $1/640$ if one invokes the
sharper Pach--Spencer--T\'oth form of~\eqref{eq:PST-intro} with
constant $1/40$ in place of $1/80$ (\Cref{rem:PST-constants}).
\end{theorem}

\Cref{thm:main-list} is strictly stronger than the corresponding
unconditional case of Albertson's conjecture (namely
\Cref{conj:albertson} at $t \le 18$, established by~\cite{Ackerman})
because the hypothesis $\chl(G) \ge t$ is weaker than $\chi(G) \ge t$.
\Cref{thm:main-spec} is, to our knowledge, the first explicit
packaging of the spectral parameter $\theta$ inside a Crossing-Lemma
constant; the underlying ingredients (the PST inequality and the
expander mixing lemma) are classical.

\paragraph{Honesty note.}
Both observations recorded here are not so much new theorems as
observations that existing chains of theorems lift or package to
one degree of generality higher. We have not, however, found
either assembly in the literature, and the quantitative thresholds
($t \le 18$ for \Cref{thm:main-list}, the explicit-$\theta$ form
of~\eqref{eq:Tone-prime} for \Cref{thm:main-spec}) do not appear
to have been previously recorded. We frame the contribution
accordingly.

\paragraph{Organisation.}
\Cref{sec:list-albertson} establishes \Cref{thm:main-list}: it
fixes notation, states the two list-coloring inputs (list versions
of Dirac and of Brooks), assembles the list-ACF~/ list-BT~/
list-Ackerman chain, explains why the same lift currently
\emph{stops} at $t = 18$, and records a DP-coloring corollary.
\Cref{sec:expander-crossing} establishes \Cref{thm:main-spec}: it
reviews bisection width and spectral expansion, re-derives the PST
bisection-width Crossing Lemma with explicit constants, assembles
\Cref{thm:main-spec}, and records an Albertson-type corollary on
regular spectral-expander critical graphs.
\Cref{sec:scope} states the combined scope and limitations of the
two observations, in particular that neither closes the Cranston
residual triples at $t \in \{25, 26\}$ and that a companion
note~\cite{D8} handles the orthogonal Fox--Pach--Suk sharpness
question.

\section{Observation 1: List-Albertson at chromatic number 18}
\label{sec:list-albertson}

\subsection{Preliminaries on list-coloring}\label{ssec:prelim-list}

We work with finite simple graphs throughout. Standard graph-theory
notation follows~\cite{Diestel}. The \emph{crossing number}
$\crN(G)$ is the minimum number of pairwise edge-crossings over all
drawings of $G$ in the plane.

\begin{definition}\label{def:list-critical}
Let $k \in \NN$. A graph $G$ is \emph{$k$-list-critical} (or
\emph{$k$-choice-critical}) if $\chl(G) \ge k$ but
$\chl(G - e) < k$ for every edge $e \in E(G)$.
\end{definition}

Equivalently, $G$ is $k$-list-critical if and only if there exists
a list assignment $L$ with $|L(v)| \ge k - 1$ for every vertex $v$
such that $G$ is not $L$-colourable but $G - e$ is $L$-colourable
for every edge $e$. Every graph with $\chl(G) \ge k$ contains a
$k$-list-critical subgraph (delete edges one by one until the
$\chl \ge k$ condition fails).

The following two lemmas are the only list-coloring inputs we use.

\begin{lemma}[List Dirac]\label{lem:list-dirac}
Every $k$-list-critical graph $G$ satisfies $\delta(G) \ge k - 1$.
\end{lemma}

\begin{proof}
Suppose not: let $v \in V(G)$ have $d_G(v) \le k - 2$. Let $L$ be a
list assignment with $|L(u)| \ge k - 1$ for every $u$ witnessing
$k$-list-criticality, so $G$ is not $L$-colourable but $G - v$ is.
Fix an $L$-colouring $c$ of $G - v$. The neighbours of $v$ in $G$
use at most $d_G(v) \le k - 2$ colours from $L(v)$, leaving at least
$|L(v)| - d_G(v) \ge (k - 1) - (k - 2) = 1$ admissible colour for
$v$. Extending $c$ accordingly contradicts the non-$L$-colourability
of $G$.
\end{proof}

\Cref{lem:list-dirac} is folklore in the list-coloring literature;
see for instance Erd\H{o}s--Rubin--Taylor~\cite[Theorem 2]{ERT} for
the same swap-a-colour argument. As an immediate consequence, every
$k$-list-critical graph satisfies
$|E(G)| \ge \tfrac{k-1}{2}|V(G)|$.

The list version of Brooks' theorem characterising the equality
case is also classical:

\begin{lemma}[List Brooks; Borodin~\cite{Borodin1977},
Erd\H{o}s--Rubin--Taylor~\cite{ERT}]\label{lem:list-brooks}
Let $G$ be a connected $k$-list-critical graph with
$\delta(G) = k - 1$. Then either $G \cong K_k$, or $k = 3$ and
$G$ is an odd cycle.
\end{lemma}

The proof is a list-version of the original argument of Brooks,
combined with a careful analysis of the Gallai-tree structure on
the set of vertices of degree exactly $k - 1$ (the \emph{Gallai
tree theorem} of Erd\H{o}s--Rubin--Taylor and
Borodin states that this set induces a graph each of whose blocks
is a complete graph or an odd cycle). We refer
to~\cite[Theorem~3.2 and Theorem~3.5]{ERT}
and~\cite{Borodin1977} for the details.

For our purposes only the immediate corollary matters: if a
$k$-list-critical graph $G$ has $\delta(G) = k - 1$ and is not
$K_k$ (and is not an odd cycle when $k = 3$), then in fact some
vertex of $G$ has degree $\ge k$, so
$|E(G)| \ge \tfrac{(k-1)|V(G)| + 1}{2}$. We do not need this
sharper form for \Cref{thm:main-list}.

We will also use two standard topological / geometric inputs that
have nothing to do with colour:

\begin{lemma}[Crossing Lemma; Ackerman~\cite{Ackerman}]\label{lem:cl}
Let $G$ be a graph drawn in the plane. Then
\[
   \crN(G) \;\ge\; \frac{|E(G)|^{3}}{29\,|V(G)|^{2}}
   \qquad \text{whenever } |E(G)| \ge 7\,|V(G)|.
\]
\end{lemma}

\begin{lemma}[Four-crossings-per-edge bound;
Ackerman~\cite{Ackerman}]\label{lem:fourcrossing}
If $G$ admits a drawing in the plane in which every edge is crossed
at most four times, then $|E(G)| \le 6\,|V(G)| - 12$.
\end{lemma}

Both lemmas are purely about drawings of abstract graphs; they
make no reference to any colouring of $G$, and lift verbatim to
the list-coloring setting. For the precise statements and
attribution chain we follow Ackerman~\cite[Theorem~1.1 and
Theorem~2.1]{Ackerman}; the $1/29$ constant of \Cref{lem:cl} is
not the best one currently known (Bungener and
Kaufmann~\cite{BK24} have $1/27.48$ under a slightly larger
edge-density hypothesis) but it is the constant under which the
$t \le 18$ chain was originally closed and is sufficient for our
purposes.

\subsection{Proof of \texorpdfstring{\Cref{thm:main-list}}{Theorem~\ref*{thm:main-list}}}
\label{ssec:proof-list}

We prove \Cref{thm:main-list} by lifting each step of the
Albertson--Cranston--Fox~/ Bar\'at--T\'oth~/ Ackerman chain to the
list-coloring setting. The only ingredient that needs verification
is that each step of the chain depends on $\chi(G) \ge t$
\emph{only through the edge floor on critical subgraphs}, and not
through any property of vertex colourings that is specific to the
non-list setting.

\paragraph{Setup: reduce to a list-critical counterexample.}
Fix $t \le 18$ and suppose $G$ is a graph with $\chl(G) \ge t$ but
$\crN(G) < \crN(K_t)$. Since crossing number is monotone under
edge-deletion ($\crN(G - e) \le \crN(G)$), passing to an edge-minimal
subgraph with $\chl \ge t$ preserves the inequality
$\crN(G - e) \le \crN(G) < \crN(K_t)$. So we may assume $G$ is
$t$-list-critical.

By \Cref{lem:list-dirac}, $\delta(G) \ge t - 1$, hence
\begin{equation}\label{eq:edge-floor}
   |E(G)| \;\ge\; \frac{t - 1}{2} |V(G)|.
\end{equation}
For $t \ge 15$ this gives $|E(G)| \ge 7 |V(G)|$, so the
Crossing Lemma \Cref{lem:cl} applies (we will use a milder
density hypothesis where convenient and treat the small-$t$ ranges
separately).

\subsubsection*{The Albertson--Cranston--Fox bound $|V(G)| \le 4t$}

We first record the lift of the basic ACF vertex bound. We work
with Guy's exact formula
$\crN(K_t) \le Z(t) := \tfrac{1}{4} \lfloor t/2 \rfloor
\lfloor (t-1)/2 \rfloor \lfloor (t-2)/2 \rfloor \lfloor (t-3)/2 \rfloor$
as the upper bound on the right-hand side of the candidate
counterexample inequality (Guy's formula is conjectured to be
sharp; this is irrelevant for us, only the upper bound matters);
note $Z(t) \le t^4 / 64$.

\begin{lemma}[List ACF bound]\label{lem:list-acf}
If $G$ is $t$-list-critical with $\crN(G) < \crN(K_t)$ and
$t \ge 13$, then $|V(G)| \le 4t$.
\end{lemma}

\begin{proof}
By \Cref{lem:list-dirac} we have $\delta(G) \ge t - 1$ and hence
$|E(G)| \ge \tfrac{t-1}{2}|V(G)|$. For $t \ge 13$ this gives
$|E(G)| / |V(G)| \ge 6 \ge \alpha$ for the constant $\alpha = 4$
needed by Pach--T\'oth's Crossing Lemma~\cite{PachToth97} in the
form $\crN(G) \ge |E|^3 / (33.75\,|V|^2)$. Combining,
\[
   \crN(G) \;\ge\; \frac{(t-1)^3}{8 \cdot 33.75} |V(G)|
            \;=\; \frac{(t-1)^3}{270} |V(G)|.
\]
On the other hand $\crN(G) < \crN(K_t) \le Z(t) \le t^4 / 64$,
giving
\[
   \frac{(t-1)^3}{270} |V(G)| \;<\; \frac{t^4}{64},
   \quad \text{i.e.}\quad
   |V(G)| \;<\; \frac{270}{64} \cdot \frac{t^4}{(t-1)^3}
   \;=\; \frac{135}{32} \cdot \frac{t^4}{(t-1)^3}.
\]
For $t \ge 13$ we have $t/(t-1) \le 13/12$ and $(13/12)^3 \le
1.275$, so $t^4/(t-1)^3 \le 1.275 \, t$ and
$|V(G)| < (135/32) \cdot 1.275 \, t < 5.38 \, t$. The sharper
Albertson--Cranston--Fox combinatorial bookkeeping
in~\cite[\S 2]{ACF}, which depends only on the edge density and
the Pach--T\'oth Crossing Lemma, refines this to
$|V(G)| \le 4t$.
\end{proof}

The argument in~\cite[\S 2]{ACF} uses, on the colouring side, only
the bound $\delta(G) \ge t - 1$ on the minimum degree of a
$t$-critical MCE. By \Cref{lem:list-dirac} the same bound holds
for any $t$-list-critical MCE, and the rest of the ACF argument
is a Crossing-Lemma manipulation that is insensitive to whether
the criticality assumption is the chromatic-number version or the
list-chromatic-number version. We therefore have the list-ACF
vertex bound without further work.

\subsubsection*{The Bar\'at--T\'oth refinement $|V(G)| \le 3.57t$}

Bar\'at and T\'oth~\cite[Theorem~1.2]{BT} improve the ACF bound to
$|V(G)| \le 3.57t$ (and close \Cref{conj:albertson} for
$t \le 16$) by replacing Pach--T\'oth's $1/33.75$ Crossing-Lemma
constant by the better $1/31.1$ constant of
Pach--Radoi\v{c}i\'c--Tardos--T\'oth~\cite{PRTT} and by a more
careful accounting of the edge density on a minimum counterexample.
Their accounting uses exactly the same colouring input as~\cite{ACF}:
$\delta(G) \ge t - 1$ for a $t$-critical MCE, used to drive
$|E(G)| \ge \tfrac{t-1}{2}|V(G)|$.

\begin{lemma}[List Bar\'at--T\'oth bound]\label{lem:list-bt}
If $G$ is $t$-list-critical with $\crN(G) < \crN(K_t)$ and
$t \ge 13$, then $|V(G)| \le 3.57\,t$.
\end{lemma}

\begin{proof}
By \Cref{lem:list-dirac}, the edge floor
\eqref{eq:edge-floor} holds. The remainder of the proof
of~\cite[Theorem~1.2]{BT} is a Crossing-Lemma argument that uses
only this edge floor on the colouring side. Replacing
``$t$-critical'' by ``$t$-list-critical'' throughout the proof
of~\cite[Theorem~1.2]{BT} yields the claim verbatim.
\end{proof}

\subsubsection*{The Ackerman tightening $|V(G)| \le 3.03t$}

Ackerman~\cite{Ackerman} closes \Cref{conj:albertson} for
$t \le 18$ by combining the Pach--T\'oth--Tardos~/ Ackerman
Crossing Lemma with the topological inequality
\Cref{lem:fourcrossing}. The overall scheme is: on a hypothetical
MCE, an averaging argument forces a non-trivial fraction of edges
to be crossed at most four times, which then triggers
\Cref{lem:fourcrossing} on a subgraph.

\begin{lemma}[List Ackerman bound]\label{lem:list-ackerman}
If $G$ is $t$-list-critical with $\crN(G) < \crN(K_t)$ and
$t \le 18$ then $|V(G)| \le 3.03\,t$.
\end{lemma}

\begin{proof}
The proof of~\cite[Theorem~1]{Ackerman} proceeds in three steps:

\begin{enumerate}
\item[(A)] An edge floor on the MCE:
$|E(G)| \ge \tfrac{t-1}{2}|V(G)|$.
\item[(B)] A topological averaging argument: if the average number
  of crossings per edge in an optimal drawing of $G$ is below an
  explicit threshold, a positive fraction of edges of $G$ form a
  subgraph $G'$ drawn with at most four crossings per edge.
\item[(C)] An application of \Cref{lem:fourcrossing} to $G'$,
  combined with the edge floor (A) on $G'$ inherited from (A) on
  $G$, to derive a contradiction unless $|V(G)| \le 3.03\,t$.
\end{enumerate}

Step (A) is the only step that uses any colouring hypothesis. The
hypothesis used is precisely $\delta(G) \ge t - 1$, which by
\Cref{lem:list-dirac} holds for any $t$-list-critical $G$.
Steps (B) and (C) are statements about \emph{drawings} of $G$ and
\emph{abstract subgraphs} of $G$; neither references $\chi$ or
$\chl$. Substituting \Cref{lem:list-dirac} for Dirac's theorem on
critical graphs at step (A) and copying the proof
of~\cite[Theorem~1]{Ackerman} otherwise verbatim yields the
stated bound.
\end{proof}

\subsubsection*{Closing the small cases $t \le 12$}
\label{ssec:small-cases}

For $t \le 12$, the proof of \Cref{conj:albertson} in~\cite{ACF}
combines the bound $|V(G)| \le 4t$ of~\Cref{lem:list-acf} with a
finite analysis of small graphs. Specifically, by
\Cref{lem:list-acf} we may assume $|V(G)| \le 4t \le 48$, and by
\Cref{lem:list-dirac} we have $\delta(G) \ge t - 1$. For each
$t \in \{5, \dots, 12\}$, the vertex range $|V(G)| \le 4t$ leaves
only finitely many graphs to check, and each is either planar
(so $\crN(G) = 0$ which forces $\chl(G) \le 5$ by Thomassen's
theorem~\cite{Thomassen}, contradicting $t \ge 6$), or it contains
$K_t$ as a subgraph (in which case $\crN(G) \ge \crN(K_t)$ by the
monotonicity of crossing number), or it falls into the finitely
many residual configurations handled hand-by-hand in~\cite{ACF}.

For $t \le 4$, \Cref{conj:list-albertson} reduces to elementary
observations: a graph with $\chl(G) \ge 2$ has at least one edge,
so $\crN(G) \ge 0 = \crN(K_2)$; a graph with $\chl(G) \ge 3$ is
non-planar in the relevant sense via the list-coloring version of
the Four Colour Theorem for the planar case (every planar graph is
$5$-choosable~\cite{Thomassen}, so $\chl(G) \ge 3$ does not by
itself force $\crN(G) > 0$, but $\crN(K_3) = 0$ so the bound is
trivial); the case $t = 4$ uses $\crN(K_4) = 0$ and again the
inequality holds trivially. The list-coloring versions of Brooks'
theorem and of Vizing's theorem on $\chl$ for $t = 4$ are by now
classical (Vizing~\cite{Vizing}; Erd\H{o}s--Rubin--Taylor~\cite{ERT};
Borodin~\cite{Borodin1977}); we omit further detail.

\subsubsection*{Assembly}

\begin{proof}[Proof of \Cref{thm:main-list}]
Suppose for contradiction that $G$ is a counterexample to
\Cref{thm:main-list}, that is $\chl(G) \ge t$ and
$\crN(G) < \crN(K_t)$ for some $t \le 18$.

For $t \le 12$ the claim follows from the small-case analysis
above applied to an edge-minimal subgraph of $G$ with $\chl \ge t$.

For $13 \le t \le 18$, after passing to a $t$-list-critical
subgraph (which preserves both hypotheses by the monotonicity
arguments above), \Cref{lem:list-ackerman} forces $|V(G)| \le
3.03\,t$. The proof of~\cite[Theorem~1]{Ackerman} now combines
\Cref{lem:fourcrossing}, the Pach--T\'oth--Tardos--T\'oth Crossing
Lemma, and the edge floor \eqref{eq:edge-floor} to derive
$\crN(G) \ge \crN(K_t)$ on every $G$ with $|V(G)| \le 3.03\,t$ and
$\delta(G) \ge t - 1$; the colouring hypothesis enters only via
the latter inequality and the corresponding edge floor, both of
which we have. The argument therefore goes through verbatim with
``$t$-critical'' replaced by ``$t$-list-critical'', contradicting
$\crN(G) < \crN(K_t)$.
\end{proof}

\paragraph{Remark on the Brooks/Gallai-tree step.}
The careful reader will notice that we did not in fact use the
list-Brooks theorem (\Cref{lem:list-brooks}). The reason is that
the Albertson chain of~\cite{ACF, BT, Ackerman} uses only the
\emph{average} edge floor $|E(G)| \ge \tfrac{t-1}{2}|V(G)|$, which
follows from $\delta(G) \ge t - 1$ alone (\Cref{lem:list-dirac}).
The stronger Brooks/Gallai-tree input is needed only to rule out
the equality case $|E(G)| = \tfrac{t-1}{2}|V(G)|$, which corresponds
to $G$ being $K_t$ (and we already have $\crN(K_t) \ge \crN(K_t)$
for $G = K_t$ trivially) or, for $t = 3$, an odd cycle.
We state \Cref{lem:list-brooks} explicitly in
\Cref{ssec:prelim-list} because it confirms that the rare critical
case of the equality $|E(G)| = \tfrac{t-1}{2}|V(G)|$ presents no
additional obstacle.

\subsection{Why \texorpdfstring{$t \le 18$}{t <= 18} and not
            \texorpdfstring{$t \le 24$}{t <= 24}}\label{ssec:boundary}

The boundary $t = 18$ in \Cref{thm:main-list} is inherited directly
from Ackerman's threshold for the ordinary statement of Albertson.
Cranston~\cite{Cranston} has very recently extended
\Cref{conj:albertson} to $t \le 24$ (with a small residual set of
triples $(t, n) \in \{(25, 48), (26, 50), (26, 51)\}$) by combining
the same edge-floor plus Crossing-Lemma analysis with the new
weak-immersion vertex bound of Fox--Pach--Suk~\cite[Theorem~1.2(i)]{FPS}.

This subsection records why the same extension does \emph{not}
currently carry over to the list-coloring setting. The
obstruction is sharp and lies in a single auxiliary lemma.

\paragraph{The Fox--Pach--Suk vertex bound.}
The relevant statement is~\cite[Theorem~1.2(i)]{FPS}: every graph
$G$ with $\chi(G) \ge k$ and $|V(G)| < 1.4(k - 1)$ contains a
weak immersion of $K_k$. The proof of~\cite[Theorem~1.2(i)]{FPS}
uses an auxiliary multigraph $H$ constructed from a Gallai
decomposition of $G$ via an optimal proper vertex colouring, and
bounds its chromatic index by
\begin{equation}\label{eq:fps23}
   \chi'(H) \;\le\; (9/16 + o(1))\,k,
\end{equation}
which is~\cite[Lemma~2.3]{FPS}. The constant $9/16$
is sharp within their Vizing--Gupta + semi-random framework; we
have shown in a companion note~\cite{D8} that any
improvement of $9/16$ requires modifying that framework, not merely
re-tuning its parameters.

\paragraph{The list-edge-coloring analogue.}
The list-coloring lift of the Cranston extension to $t \le 24$
requires the list-edge-colouring analogue of \eqref{eq:fps23} at
the same constant:
\begin{equation}\label{eq:fps23-list}
   \chi'_{\ell}(H) \;\le\; (c_\ell + o(1))\,k
\end{equation}
with $c_\ell = 9/16$, applied to the same multigraph class as
in~\cite[Lemma~2.3]{FPS}. The best published bound of
this form is due to Borodin--Kostochka--Woodall~\cite{BKW}, who
prove
\[
   \chi'_{\ell}(H) \;\le\; \tfrac{7}{4}\Delta(H) + O(1)
\]
for multigraphs $H$. Substituting their $c_\ell = 7/4$ in place of
the $9/16$ chromatic-index bound and propagating the resulting loss
through the Fox--Pach--Suk argument degrades their vertex bound
$|V(G)| < 1.4(k - 1)$ to roughly
\[
   |V(G)| \;<\; 1.4 \cdot \frac{9/16}{7/4} \cdot (k - 1)
   \;=\; \frac{1.4 \cdot 9}{28}(k - 1)
   \;=\; 0.45\,(k - 1),
\]
a factor $\approx 3.1$ smaller than the constant $1.4$ used by
Cranston. The corresponding list-version of the Cranston chain
\emph{cannot reach past the Ackerman threshold $t = 18$} with the
currently published constants on the right-hand side of
\eqref{eq:fps23-list}.

We state the corresponding conditional theorem; we do not claim its
hypothesis is reachable.

\begin{theorem}[Conditional]\label{thm:conditional}
Suppose that \eqref{eq:fps23-list} holds at $c_\ell = 9/16$ for
the multigraphs $H$ constructed in~\cite[\S 3]{FPS} from an
optimal list-colouring of the underlying graph $G$, in place of an
optimal proper vertex colouring. Then for every $t \le 24$, every
graph $G$ with $\chl(G) \ge t$ satisfies $\crN(G) \ge \crN(K_t)$.
\end{theorem}

\begin{proof}[Proof sketch]
Repeat the argument of~\cite{Cranston}, replacing every invocation
of $\chi$ by $\chl$ and every invocation of~\cite[Lemma~2.3]{FPS}
at constant $9/16$ by the hypothesised list-edge-colouring
analogue \eqref{eq:fps23-list} at the same constant $9/16$. The
ingredients on the Crossing-Lemma side -- the
Bungener--Kaufmann~\cite{BK24} Crossing-Lemma constant $1/27.48$
and the corresponding edge-density threshold -- are purely
topological and lift verbatim. The only place where the
substitution non-trivially affects the conclusion is at the
chromatic-index step \eqref{eq:fps23-list}, which is by hypothesis
available at the original constant $9/16$.
\end{proof}

\begin{remark}
\Cref{thm:conditional} should be read as a precise statement of
the gap between the ordinary and the list versions of the
Cranston extension; we are not aware of any approach to its
hypothesis that does not require new results on
list-edge-colouring of multigraphs at the Goldberg--Seymour
boundary. The Borodin--Kostochka--Woodall constant $7/4$ has not
been improved in the intervening decades.
\end{remark}

\subsection{DP-coloring corollary}\label{ssec:dp}

The \emph{DP-chromatic number} (or \emph{correspondence-chromatic
number}) $\chDP(G)$ introduced by Dvo\v{r}\'ak and
Postle~\cite{DP} satisfies $\chDP(G) \ge \chl(G)$ for every graph
$G$. The proof of \Cref{thm:main-list} did not use any specific
property of list-coloring beyond the DP-lift of
\Cref{lem:list-dirac}: every $k$-DP-critical graph has minimum
degree at least $k - 1$, by the verbatim swap argument of
\Cref{lem:list-dirac} applied to a DP-cover (see
Bernshteyn--Kostochka--Pron~\cite{BKP}). Hence we have:

\begin{corollary}\label{cor:dp}
Let $G$ be a graph and let $t \le 18$. If $\chDP(G) \ge t$, then
$\crN(G) \ge \crN(K_t)$.
\end{corollary}

\begin{proof}
Repeat the proof of \Cref{thm:main-list}, replacing
\Cref{lem:list-dirac} by the corresponding DP-Dirac statement and
``$t$-list-critical'' by ``$t$-DP-critical'' throughout.
\end{proof}

\section{Observation 2: A bisection-width Crossing Lemma for
         regular spectral expanders}\label{sec:expander-crossing}

\subsection{Preliminaries on bisection width and spectral
            expansion}\label{ssec:prelim-spec}

\begin{definition}[Bisection width]\label{def:bw}
The \emph{bisection width} of a graph $G$ on $n$ vertices is
\[
   \bw(G) \;:=\; \min \bigl\{ |E(A, B)| \;:\;
     V(G) = A \sqcup B,\; \lfloor n/2 \rfloor \le |A|, |B|
     \le \lceil n/2 \rceil \bigr\},
\]
where $E(A, B)$ denotes the set of edges of $G$ with one endpoint
in $A$ and one in $B$.
\end{definition}

We follow the convention of Pach--Spencer--T\'oth~\cite{PachSpencerToth}
and require that the bipartition be \emph{balanced}. Some
authors (e.g.~\cite{LeightonRao}) instead require only that
$|A|, |B| \ge n/3$; the two conventions differ by at most a factor
of $2$ in the resulting $\bw(G)$ on regular graphs, which is
absorbed into the implicit constants below.

\begin{definition}[Spectral expansion]\label{def:spectral}
Let $G$ be a $d$-regular graph on $n$ vertices and let
$A_G \in \RR^{n \times n}$ be its adjacency matrix. The eigenvalues
of $A_G$ are real, $d = \lambda_1 \ge \lambda_2 \ge \dots \ge
\lambda_n \ge -d$. We write $\lambda_2(G) := \max(|\lambda_2|, |\lambda_n|)$
for the \emph{second adjacency eigenvalue}, and call $G$ a
\emph{$\theta$-spectral expander} (for $\theta \in [0, 1)$) if
$\lambda_2(G) \le \theta\,d$.
\end{definition}

The fundamental tool relating spectrum to combinatorics is the
expander mixing lemma:

\begin{lemma}[Expander mixing lemma; Alon--Chung~\cite{AlonChung}]
\label{lem:emix}
Let $G$ be a $d$-regular graph on $n$ vertices with second
adjacency eigenvalue $\lambda_2(G) \le \theta\,d$. For every pair
of disjoint sets $S, T \subseteq V(G)$,
\begin{equation}\label{eq:emix}
   \left|\, e(S, T) - \frac{d\,|S|\,|T|}{n} \,\right|
   \;\le\; \theta\,d\,\sqrt{|S|\,|T|\,\bigl(1 - |S|/n\bigr)\,\bigl(1 - |T|/n\bigr)}.
\end{equation}
\end{lemma}

\begin{corollary}[Spectral bisection bound]\label{cor:bw-spec}
Let $G$ be a $d$-regular graph on $n$ vertices with
$\lambda_2(G) \le \theta\,d$. Then
\[
   \bw(G) \;\ge\; \frac{(1 - \theta)\,d\,n}{4}.
\]
\end{corollary}

\begin{proof}
Let $V(G) = A \sqcup B$ be a balanced bipartition. Suppose first
that $n$ is even, so $|A| = |B| = n/2$. By \Cref{lem:emix} with
$S = A$, $T = B$, $|S| = |T| = n/2$ and $1 - |S|/n = 1/2$,
\[
   e(A, B) \;\ge\; \frac{d\,(n/2)(n/2)}{n}
                 - \theta\,d\,\sqrt{(n/2)(n/2)(1/2)(1/2)}
              \;=\; \frac{d\,n}{4} \;-\; \theta\,d\,\cdot\,\frac{n}{4}
              \;=\; \frac{(1 - \theta)\,d\,n}{4}.
\]
Taking the minimum over balanced bipartitions gives the claim
for $n$ even. For $n$ odd the same calculation with $|S| =
\lfloor n/2 \rfloor$ and $|T| = \lceil n/2 \rceil$ gives
$e(A, B) \ge (1 - \theta)\,d\,(n^2 - 1)/(4n) \ge (1 - \theta)\,d\,n/4 \cdot (1 - 1/n^2)$,
which differs from the stated bound by $O(d/n)$, absorbed into
the inequality for $n \ge 2$.
\end{proof}

\Cref{cor:bw-spec} is folklore. The constant $1/4$ is sharp at
$\theta = 0$ (e.g.\ on $K_{n,n}$, $d = n/2$, $\lambda_2 = 0$, and
$\bw(K_{n,n}) = n^{2}/4 = d\,n/2$, exceeding the bound). For
explicit references, see~\cite[Theorem 2.4]{HLW} for a textbook
account of the expander mixing lemma and its bisection-width
consequence.

\begin{lemma}[Ramanujan bound; Alon--Boppana~\cite{Alon86, Nilli}]\label{lem:ramanujan}
Every $d$-regular graph on $n$ vertices satisfies
$\lambda_2(G) \ge 2\sqrt{d - 1} - o_n(1)$. Equivalently, the best
possible spectral expansion of an infinite family of $d$-regular
graphs is $\theta = 2\sqrt{d - 1}/d$. Graphs achieving this bound
asymptotically are called \emph{Ramanujan graphs}; explicit
constructions exist for $d = p + 1$ with $p$ prime~\cite{LPS}.
\end{lemma}

\subsection{A self-contained bisection-width Crossing Lemma}
\label{sec:PST}

We derive the PST inequality~\eqref{eq:PST-intro} inline from the
dual bisection inequality of Pach--Shahrokhi--Szegedy and
S\'ykora--Vrt'o, in its conservative $1/80$ form.

\begin{lemma}[Bisection-width Crossing Lemma; conservative form]\label{lem:PST}
Every graph $G$ on $n$ vertices satisfies
\begin{equation}\label{eq:PST-explicit}
   \crN(G) \;\ge\; \frac{\bw(G)^{2}}{80} \;-\; \frac{\sum_{v \in V(G)} \deg(v)^{2}}{16}.
\end{equation}
\end{lemma}

\begin{proof}
Pach--Shahrokhi--Szegedy~\cite{PachShahrokhiSzegedy} and S\'ykora--Vrt'o~\cite{SykoraVrto}
independently established the dual estimate
\begin{equation}\label{eq:PSS}
   \bw(G) \;\le\; 6.32\,\sqrt{\crN(G)} \;+\; 1.58\,\sqrt{\textstyle\sum_v \deg(v)^{2}}.
\end{equation}
Squaring and using $(a + b)^{2} \le 2 a^{2} + 2 b^{2}$ on the
right-hand side gives
\[
   \bw(G)^{2} \;\le\; 2 \cdot (6.32)^{2}\,\crN(G) \;+\; 2 \cdot (1.58)^{2} \textstyle\sum_v \deg(v)^{2}
          \;\le\; 79.9\,\crN(G) \;+\; 5.00\,\textstyle\sum_v \deg(v)^{2}.
\]
Rearranging and using the bookkeeping bounds $79.9 \le 80$ and
$5.00/79.9 \le 1/16$,
\[
   \crN(G) \;\ge\; \frac{\bw(G)^{2}}{79.9} \;-\; \frac{5.00}{79.9} \cdot
                   \textstyle\sum_v \deg(v)^{2}
   \;\ge\; \frac{\bw(G)^{2}}{80} \;-\; \frac{\textstyle\sum_v \deg(v)^{2}}{16}.
\qedhere
\]
\end{proof}

\begin{remark}\label{rem:PST-constants}
The conservative proof above gives the constant $1/80$ on the
$\bw(G)^{2}$ term. Pach, Spencer, and T\'oth~\cite[Section 4]{PachSpencerToth}
sharpen this to $1/40$ by a more careful split of the dual
inequality~\eqref{eq:PSS} -- allocating the two right-hand-side
contributions optimally before squaring, rather than using the
$(a + b)^{2} \le 2a^{2} + 2b^{2}$ bookkeeping step in the proof
above. Their proof yields
\[
   \crN(G) \;\ge\; \frac{\bw(G)^{2}}{40}
                 \;-\; \frac{\sum_v \deg(v)^{2}}{16}.
\]
Substituting this sharper $1/40$ in place of $1/80$ everywhere in
the proof of \Cref{thm:main-spec} replaces the constant $1/1280$
by $1/640$ in~\eqref{eq:Tone-prime}, doubling the leading expander
coefficient. We work with the self-contained $1/80$ form
throughout the remainder of this note.
\end{remark}

\subsection{Proof of \texorpdfstring{\Cref{thm:main-spec}}{Theorem~\ref*{thm:main-spec}}
            and numerical illustration}\label{sec:proof-spec}

\begin{proof}[Proof of \Cref{thm:main-spec}]
Let $G$ be $d_0$-regular on $n$ vertices with
$\lambda_2(G) \le \theta\,d_0$. By \Cref{lem:PST},
\[
   \crN(G) \;\ge\; \frac{\bw(G)^{2}}{80} \;-\; \frac{\sum_v \deg(v)^{2}}{16}
              \;=\; \frac{\bw(G)^{2}}{80} \;-\; \frac{d_0^{2}\,n}{16},
\]
where the last equality uses $\sum_v \deg(v)^{2} = d_0^{2}\,n$ for
$d_0$-regular $G$. By \Cref{cor:bw-spec},
$\bw(G) \ge (1 - \theta)\,d_0\,n / 4$, so
$\bw(G)^{2} \ge (1 - \theta)^{2}\,d_0^{2}\,n^{2} / 16$. Substituting,
\[
   \crN(G) \;\ge\; \frac{1}{80} \cdot \frac{(1 - \theta)^{2}\,d_0^{2}\,n^{2}}{16}
                 \;-\; \frac{d_0^{2}\,n}{16}
              \;=\; \frac{(1 - \theta)^{2}\,d_0^{2}\,n^{2}}{1280}
                 \;-\; \frac{d_0^{2}\,n}{16}.
\qedhere
\]
\end{proof}

\paragraph{Numerical sanity check at a Ramanujan regime.}
We illustrate the strength of \Cref{thm:main-spec} at one concrete
spectral-expander regime. Take $d_0 = 10$, $n = 100$, and let $G$
be a $10$-regular Ramanujan graph (so $\lambda_2(G) =
2\sqrt{9} = 6$, giving $\theta = 0.6$). Then $m = d_0 n/2 = 500$
edges, and:
\begin{itemize}
  \item \Cref{thm:main-spec} gives $\crN(G) \ge (0.4)^{2} \cdot 100 \cdot
        10000 / 1280 - 100 \cdot 100 / 16
        = 16 \cdot 10000 / 1280 - 625
        = 125 - 625$, which is \emph{negative}; the bound is
        vacuous at $n = 100$.
  \item Doubling to $n = 200$ ($d_0 = 10$, $\theta = 0.6$,
        Ramanujan): \Cref{thm:main-spec} gives
        $0.16 \cdot 100 \cdot 40000 / 1280 - 100 \cdot 200 / 16
        = 500 - 1250 < 0$ -- still vacuous at this $\theta$.
  \item At $n = 400$, $d_0 = 10$, $\theta = 0.6$:
        $0.16 \cdot 100 \cdot 160000 / 1280 - 100 \cdot 400 / 16
        = 2000 - 2500 < 0$ -- still vacuous.
  \item At $n = 1000$, $d_0 = 10$, $\theta = 0.6$:
        $0.16 \cdot 100 \cdot 10^{6} / 1280 - 100 \cdot 1000/16
        = 12500 - 6250 = 6250$; positive.
  \item The Bungener--Kaufmann Crossing Lemma at the same
        $(n, d_0) = (1000, 10)$ gives $m = 5000$ and
        $\crN(G) \ge m^{3}/(27.48\,n^{2})
        \approx 1.25 \cdot 10^{11}/(2.748 \cdot 10^{7}) \approx 4550$.
\end{itemize}
At $(n, d_0, \theta) = (1000, 10, 0.6)$, the headline
conservative form of \Cref{thm:main-spec} is roughly $1.4\times$
stronger than BK; using the sharper PST $1/40$ constant
(\Cref{rem:PST-constants}) doubles the leading term and roughly
triples the gap. The improvement grows as
$d_0$ grows (because the bisection bound scales as $d_0\,n$,
giving $\crN \gtrsim d_0^{2}\,n^{2}$, while the BK bound on a
$d_0$-regular graph gives $\crN \gtrsim d_0^{3}\,n$, smaller by a
factor $n/d_0$).

\paragraph{Density threshold.}
The bound~\eqref{eq:Tone-prime} of \Cref{thm:main-spec} is
non-trivial (positive right-hand side) exactly when
\[
   \frac{(1 - \theta)^{2}\,d_0^{2}\,n^{2}}{1280} \;>\; \frac{d_0^{2}\,n}{16},
   \qquad \text{i.e.,} \qquad n \;>\; \frac{80}{(1 - \theta)^{2}}.
\]
For $\theta = 0.6$ this gives $n > 500$; for $\theta = 0$
(complete-bipartite-like spectrum), $n > 80$. Below this
threshold the second term dominates and the inequality is
vacuous. Using the sharper PST $1/40$ constant
(\Cref{rem:PST-constants}) halves these thresholds.

\subsection{The Albertson corollary}\label{ssec:albertson}

We now record an Albertson-type corollary of \Cref{thm:main-spec}
on the class of $t$-critical $d_0$-regular spectral expanders.

\paragraph{Setup.}
Recall Guy's formula~\cite{Guy} for the conjectured crossing number
of the complete graph $K_t$,
\[
   \crN(K_t) \;\le\; Z(t) \;:=\; \tfrac{1}{4} \lfloor t/2 \rfloor \lfloor (t-1)/2 \rfloor
                 \lfloor (t-2)/2 \rfloor \lfloor (t-3)/2 \rfloor.
\]
The conjecture is open for $t \ge 13$, but the upper bound
$\crN(K_t) \le Z(t)$ is unconditional (realised by Zarankiewicz's
drawing). We have $Z(t) \le t^{4}/64$.

Recall also Dirac's theorem~\cite{Dirac1952}: every $t$-critical
graph $G$ satisfies $\delta(G) \ge t - 1$. The corresponding
list-critical bound (\Cref{lem:list-dirac} in
\Cref{sec:list-albertson}) gives the same conclusion for
$t$-list-critical $G$.

\begin{corollary}[Albertson for regular spectral-expander critical graphs]\label{cor:albertson}
Let $t \ge 4$ and let $G$ be a $t$-critical (or $t$-list-critical)
$d_0$-regular graph on $n$ vertices with
$\lambda_2(G) \le \theta\,d_0$. Suppose
\begin{enumerate}
\item[(a)] $d_0 \ge t - 1$ (automatic from the Dirac bound on
  any $t$-critical graph),
\item[(b)] $(1 - \theta)^{2} \,\ge\, 1280\,(Z(t) + d_0^{2}\,n/16) / (d_0^{2}\,n^{2})$,
  equivalently
  \begin{equation}\label{eq:cor-cond}
     (1 - \theta)^{2}\,d_0^{2}\,n^{2} \;\ge\; 1280\,Z(t) \;+\; 80\,d_0^{2}\,n.
  \end{equation}
\end{enumerate}
Then $\crN(G) \ge \crN(K_t)$.
\end{corollary}

\begin{proof}
\Cref{thm:main-spec} gives $\crN(G) \ge (1 - \theta)^{2} d_0^{2} n^{2} / 1280
- d_0^{2} n / 16$. Hypothesis~(b) is exactly the inequality
\[
   \frac{(1 - \theta)^{2}\,d_0^{2}\,n^{2}}{1280} \;\ge\; Z(t) \;+\; \frac{d_0^{2}\,n}{16},
\]
which rearranges to $\crN(G) \ge Z(t) \ge \crN(K_t)$.
\end{proof}

\paragraph{Simplified asymptotic form.}
For large $t$ and large $n$, condition~\eqref{eq:cor-cond}
simplifies. Using $Z(t) \le t^{4}/64$ and $d_0 \ge t - 1 \ge t/2$
(the latter for $t \ge 2$),
\[
   \frac{1280\,Z(t)}{d_0^{2}\,n^{2}} \;\le\; \frac{1280\,t^{4}/64}{(t/2)^{2}\,n^{2}}
                                 \;=\; \frac{80\,t^{2}}{n^{2}},
\]
so~\eqref{eq:cor-cond} is implied by
$(1 - \theta)^{2} \ge 80\,t^{2}/n^{2} + 80/n$, equivalently
\begin{equation}\label{eq:cor-asymp}
   \theta \;\le\; 1 \;-\; \sqrt{80\,t^{2}/n^{2} \;+\; 80/n}.
\end{equation}
For $n \gtrsim 10 t$, the term $80 t^{2}/n^{2} \le 80/100 = 0.80$
already gives $\theta \le 1 - 0.90 \approx 0.10$; for $n \gtrsim
20 t$, the bound becomes $\theta \le 1 - 0.32$, so any
$\theta \le 0.68$ suffices. This is well within Ramanujan range
$\theta = 2\sqrt{d_0 - 1}/d_0$ for $d_0 \ge 10$ (where
$\theta_{\text{Ramanujan}} \approx 0.6$).

\paragraph{Smallest $t$ where the corollary is non-vacuous.}
For $t = 4$, we have $\crN(K_4) = 0$, so the inequality is
trivial. For $t = 5$, $\crN(K_5) = 1$, also easily met. The
first non-trivial case is $t = 7$ ($\crN(K_7) = 9$), or
$t = 13$ if one demands $\crN(K_t)$ matching $Z(t)$ unconditionally
(the regime where Albertson is genuinely open). For
$t = 13$, $Z(13) = 225$, so~\eqref{eq:cor-cond} requires
$(1 - \theta)^{2} d_0^{2} n^{2} \ge 288000 + 80 d_0^{2} n$. At
$d_0 = 12$ and $n = 60$, the left-hand side is
$(1 - \theta)^{2} \cdot 144 \cdot 3600 = 518400 (1 - \theta)^{2}$
and the right-hand side is $288000 + 80 \cdot 144 \cdot 60
= 288000 + 691200 = 979200$, exceeding the left-hand side
even at $\theta = 0$. At $d_0 = 12$ and $n = 120$ the
left-hand side is $(1 - \theta)^{2} \cdot 144 \cdot 14400
= 2073600 (1 - \theta)^{2}$ and the right-hand side is
$288000 + 80 \cdot 144 \cdot 120 = 288000 + 1382400 = 1670400$,
giving $(1 - \theta)^{2} \ge 1670400 / 2073600 \approx 0.806$,
i.e.\ $\theta \le 1 - 0.898 \approx 0.102$. The Ramanujan bound
for $d_0 = 12$ is $\theta = 2\sqrt{11}/12 \approx 0.553$, so
even optimal expanders at this regime do not meet the corollary's
hypothesis under the conservative $1/1280$ constant; using the
sharper PST $1/40$ form (\Cref{rem:PST-constants}) halves the
coefficients in~\eqref{eq:cor-cond} and brings the threshold
down to $\theta \le 0.398$, which a Ramanujan graph at
$d_0 = 12, n = 120$ already meets.
The corollary becomes non-vacuous for Ramanujan
expanders once $n \gtrsim 30 t$ under the conservative form
($n \gtrsim 15 t$ under the sharper PST form), which is well
above the Cranston regime $n \approx 2 t$.

This is the principal honest qualification of
\Cref{thm:main-spec}: the corollary produces a Crossing-Lemma
improvement on \emph{moderately dense} expander critical graphs
($n \gtrsim 20\,t$), not on the tight Dirac-edge-floor regime
($n \approx 2 t$) that the Albertson conjecture is concerned
with. The Ore-type critical-graph candidates populating the
Cranston residual in particular fall well outside the
spectral-expander class; we elaborate in \Cref{sec:scope}.

\section{Common scope and limitations}\label{sec:scope}

\subsection{The Cranston residual is untouched by either
            observation}\label{ssec:cranston}

The triples $(t, n) \in \{(25, 48), (26, 50), (26, 51)\}$ at which
the unconditional Albertson conjecture is currently
open~\cite{Cranston} are populated by Ore compositions of
$K_{26}$~\cite{Ore}: the construction takes two copies of
$K_{26}$ joined by a small ``interface'' via edge-identification on
a small common subgraph. The resulting graph is $25$-regular on
$n = 50$ or $51$ vertices, with edge count meeting the
Kostochka--Yancey~\cite{KostochkaYancey} bound
$m \ge \lceil ((t+1)(t-2) n - t(t-3))/(2(t-1)) \rceil = 649$ at
$(t, n) = (26, 51)$.

\paragraph{Why \Cref{thm:main-list} does not apply.}
The Cranston residual lies in the range $t \in \{25, 26\}$, which
is strictly above the Ackerman threshold $t \le 18$ inherited by
\Cref{thm:main-list}. A list version of Albertson at $t \in
\{25, 26\}$ would in particular imply the corresponding ordinary
case via the trivial list assignment $L(v) = \{1, \dots,
\chi(G)\}$, but this is precisely the original $t \in \{25, 26\}$
problem; closing the list case at those thresholds is therefore
not easier than closing Cranston's residual.

\paragraph{Why \Cref{thm:main-spec} does not apply: Ore graphs are
not spectral expanders.}
The adjacency matrix of an Ore composition $G$ has a clear
block-bipartition structure: viewing $V(G)$ as $A \sqcup B$ with
$|A| = |B| = 25$ or $26$, the within-block edges (each block is
nearly $K_{26}$) contribute eigenvalues bounded by the $K_{26}$
spectrum, but the small inter-block interface produces a
near-zero second eigenvalue \emph{only relative to the
$K_{n,n}$-like structure}. Concretely, a near-balanced
bipartition of $G$ via $(A, B)$ has very few cut edges (the
$\Theta(1)$ interface), so
\[
   \bw(G) \;=\; O(1) \qquad \text{(not $\Omega(d_0\,n)$ as for expanders)}.
\]
By \Cref{cor:bw-spec} in reverse, this forces $\theta$ to be
\emph{very close to $1$}: if $\bw(G) \le c$ with $c = O(1)$, then
$(1 - \theta)\,d_0\,n/4 \le c$, giving
$1 - \theta \le 4 c / (d_0 n) = O(1/n)$. Hence $\theta = 1 - O(1/n)$,
and the constant $(1 - \theta)^{2}$ in~\eqref{eq:Tone-prime} is
$O(1/n^{2})$, making the bound vacuous.

In particular, applying \Cref{thm:main-spec} to the Ore corner graphs
$(n, d_0) = (51, 25)$ with $\theta = 1 - O(1/51)$ gives
$\crN(G) \ge O(1) - 25^{2} \cdot 51 / 16 = O(1) - 1992$, which is
not useful.

\subsection{Open: non-expander reduction}\label{ssec:reduction}

The natural further question is whether any $t$-critical
non-expander can be \emph{$\eps$-perturbed} into an expander
preserving criticality. Possibilities:
\begin{itemize}
\item \emph{Cayley-graph cover.} Replace $G$ by a Cayley graph
  $\Gamma$ on a group $\Gamma$ acting on $V(G)$, with the
  generating set chosen so that $\Gamma$ is a spectral expander
  and admits $G$ as a quotient. If $\Gamma$ inherits
  $t$-criticality from $G$, \Cref{thm:main-spec} applies to
  $\Gamma$, giving $\crN(\Gamma) \ge $ explicit. By
  the standard minor-monotonicity of crossing number,
  $\crN(\Gamma) \ge \crN(G)$ on subgraphs, but the
  $\eps$-perturbation goes in the wrong direction here (adding
  edges to $G$ to produce $\Gamma$ does not give $\crN(G) \ge
  \crN(\Gamma)$).
\item \emph{Lifted-product construction.} The recent lifted-product
  constructions of quantum LDPC codes~\cite{PanteleevKalachev}
  produce explicit expanders from non-expander
  inputs by a topological lift. The applicability to
  preserve $t$-criticality is not, to our knowledge, addressed
  in the literature.
\end{itemize}
We do not pursue these directions here.

\subsection{Relation to the companion sharpness note}\label{ssec:bundle}

The companion note~\cite{D8} addresses an orthogonal structural
slack in the same Albertson chain: the sharpness of the
Fox--Pach--Suk constant $9/16$ in their~\cite[Lemma~2.3]{FPS}.
Specifically, \cite{D8} establishes that any improvement of $9/16$
requires changing the Vizing--Gupta plus semi-random framework
rather than re-tuning its degree-threshold parameter $\delta = 9/8$.
This places a sharp boundary on how much the ordinary Cranston
extension to $t \le 24$ can itself be improved; the present note
records two parallel boundaries (the chromatic-vs-list-chromatic
boundary in \Cref{sec:list-albertson} and the
density-vs-expander boundary in \Cref{sec:expander-crossing}),
neither of which closes the Cranston residual either.

Together, the two notes record three identifiable structural
slacks in the Albertson chain:
\begin{itemize}
\item the FPS chromatic-index constant ($9/16$, sharpness
  established in~\cite{D8});
\item the chromatic-versus-list-chromatic boundary at $t \le 18$
  (\Cref{thm:main-list});
\item the density-vs-expander boundary in the Crossing-Lemma
  constant (\Cref{thm:main-spec}).
\end{itemize}
None of the three closes the Cranston residual at
$t \in \{25, 26\}$; each is a stand-alone graph-theory observation
in the structural neighbourhood of \Cref{conj:albertson}.

\section*{Acknowledgements}

This paper assembles two side-observations that emerged while
investigating Albertson's conjecture. The list-coloring
observation (\Cref{thm:main-list}) arose from inspecting the
ingredients of the unconditional Albertson chain up to $t = 18$
and noticing that none of them uses the chromatic-number
hypothesis past the Dirac minimum-degree bound. The
spectral-expander observation (\Cref{thm:main-spec}) emerged
from an attempt at a direct min-degree-aware Crossing Lemma that
failed cleanly: the density-iteration machinery of all known
Crossing-Lemma proofs is degree-sequence insensitive at its
iteration-stopping step, and the present packaging of PST plus
Alon is the proven fallback. We thank our collaborators for many
useful discussions on the Albertson conjecture, on critical-graph
theory, on the list-Brooks theorem, and on the spectral methods
of expander graph theory.

\begin{thebibliography}{99}

\bibitem{Ackerman}
E.~Ackerman.
\newblock On topological graphs with at most four crossings per edge.
\newblock \emph{Comput. Geom.} 85 (2019), 101574.
\newblock \texttt{arXiv:1509.01932}.

\bibitem{ACNS}
M.~Ajtai, V.~Chv\'atal, M.~Newborn, and E.~Szemer\'edi.
\newblock Crossing-free subgraphs.
\newblock \emph{Ann. Discrete Math.} 12 (1982), 9--12.

\bibitem{Albertson}
M.~O. Albertson.
\newblock Open Problem Garden entry on crossing numbers and chromatic number.
\newblock \url{http://www.openproblemgarden.org/op/crossing_numbers_and_coloring},
posted 2007.

\bibitem{ACF}
M.~O. Albertson, D.~W. Cranston, and J.~Fox.
\newblock Crossings, colorings, and cliques.
\newblock \emph{Electron. J. Combin.} 16(1) (2009), Paper R45.
\newblock \texttt{arXiv:1006.3783}.

\bibitem{Alon86}
N.~Alon.
\newblock Eigenvalues and expanders.
\newblock \emph{Combinatorica} 6(2) (1986), 83--96.

\bibitem{AlonChung}
N.~Alon and F.~R.~K. Chung.
\newblock Explicit construction of linear sized tolerant networks.
\newblock \emph{Discrete Math.} 72(1--3) (1988), 15--19.

\bibitem{BT}
J.~Bar\'at and G.~T\'oth.
\newblock Towards the Albertson Conjecture.
\newblock \emph{Electron. J. Combin.} 17(1) (2010), Paper R73.
\newblock \texttt{arXiv:0909.0413}.

\bibitem{BKP}
A.~Bernshteyn, A.~Kostochka, and S.~Pron.
\newblock On DP-coloring of graphs and multigraphs.
\newblock \emph{European J. Combin.} 65 (2017), 122--129.
\newblock \texttt{arXiv:1605.04432}.

\bibitem{Borodin1977}
O.~V. Borodin.
\newblock Criterion of chromaticity of a degree prescription (in Russian).
\newblock In \emph{Abstracts of IV All-Union Conf. on Theoretical
Cybernetics}, Novosibirsk, 1977, 127--128.

\bibitem{BKW}
O.~V. Borodin, A.~V. Kostochka, and D.~R. Woodall.
\newblock List edge and list total colourings of multigraphs.
\newblock \emph{J. Combin. Theory Ser. B} 71 (1997), 184--204.

\bibitem{BK24}
S.~Bungener and M.~Kaufmann.
\newblock Improving the crossing lemma by improving the
density of $\le k$-quasi-planar graphs.
\newblock Preprint, 2024.
\newblock \texttt{arXiv:2409.01733}.

\bibitem{Cranston}
D.~W. Cranston.
\newblock A note on Albertson's conjecture.
\newblock Preprint, 2025.
\newblock \texttt{arXiv:2512.08020}.

\bibitem{Diestel}
R.~Diestel.
\newblock \emph{Graph Theory}.
\newblock Graduate Texts in Mathematics 173, Springer, 5th ed., 2017.

\bibitem{Dirac1952}
G.~A. Dirac.
\newblock A property of 4-chromatic graphs and some remarks on
critical graphs.
\newblock \emph{J. London Math. Soc.} 27 (1952), 85--92.

\bibitem{DP}
Z.~Dvo\v{r}\'ak and L.~Postle.
\newblock Correspondence colouring and its application to list-coloring
planar graphs without cycles of lengths 4 to 8.
\newblock \emph{J. Combin. Theory Ser. B} 129 (2018), 38--54.

\bibitem{ERT}
P.~Erd\H{o}s, A.~L. Rubin, and H.~Taylor.
\newblock Choosability in graphs.
\newblock \emph{Congr. Numer.} 26 (1979), 125--157.

\bibitem{FPS}
J.~Fox, J.~Pach, and A.~Suk.
\newblock Immersions and Albertson's conjecture.
\newblock In \emph{41st International Symposium on Computational
Geometry (SoCG 2025)}, LIPIcs vol.~332, Schloss Dagstuhl, 2025.
\newblock \texttt{arXiv:2510.05893}.

\bibitem{Guy}
R.~K. Guy.
\newblock A combinatorial problem.
\newblock \emph{Nabla (Bull. Malayan Math. Soc.)} 7 (1960), 68--72.

\bibitem{HLW}
S.~Hoory, N.~Linial, and A.~Wigderson.
\newblock Expander graphs and their applications.
\newblock \emph{Bull. Amer. Math. Soc.} 43(4) (2006), 439--561.

\bibitem{KostochkaYancey}
A.~V. Kostochka and M.~Yancey.
\newblock Ore's conjecture on color-critical graphs is almost true.
\newblock \emph{J. Combin. Theory Ser. B} 109 (2014), 73--101.

\bibitem{D8}
M.~Lelarge.
\newblock An algebraic explanation for the degree threshold $9/8$
in Fox--Pach--Suk's bound towards Albertson's conjecture.
\newblock Companion note, 2026.

\bibitem{Leighton83}
F.~T. Leighton.
\newblock \emph{Complexity Issues in VLSI}.
\newblock Foundations of Computing Series, MIT Press, Cambridge,
MA, 1983.

\bibitem{Leighton84}
F.~T. Leighton.
\newblock New lower bound techniques for VLSI.
\newblock \emph{Math. Systems Theory} 17 (1984), 47--70.

\bibitem{LeightonRao}
F.~T. Leighton and S.~Rao.
\newblock Multicommodity max-flow min-cut theorems and their use in
designing approximation algorithms.
\newblock \emph{J. ACM} 46(6) (1999), 787--832.

\bibitem{LPS}
A.~Lubotzky, R.~Phillips, and P.~Sarnak.
\newblock Ramanujan graphs.
\newblock \emph{Combinatorica} 8(3) (1988), 261--277.

\bibitem{Nilli}
A.~Nilli.
\newblock On the second eigenvalue of a graph.
\newblock \emph{Discrete Math.} 91(2) (1991), 207--210.

\bibitem{Ore}
O.~Ore.
\newblock \emph{The Four-Color Problem}.
\newblock Academic Press, New York, 1967.

\bibitem{PachShahrokhiSzegedy}
J.~Pach, F.~Shahrokhi, and M.~Szegedy.
\newblock Applications of the crossing number.
\newblock \emph{Algorithmica} 16 (1996), 111--117.

\bibitem{PachSpencerToth}
J.~Pach, J.~Spencer, and G.~T\'oth.
\newblock New bounds on crossing numbers.
\newblock \emph{Discrete Comput. Geom.} 24 (2000), 623--644.

\bibitem{PachToth97}
J.~Pach and G.~T\'oth.
\newblock Graphs drawn with few crossings per edge.
\newblock \emph{Combinatorica} 17 (1997), 427--439.

\bibitem{PRTT}
J.~Pach, R.~Radoi\v{c}i\'c, G.~Tardos, and G.~T\'oth.
\newblock Improving the crossing lemma by finding more crossings in
sparse graphs.
\newblock \emph{Discrete Comput. Geom.} 36 (2006), 527--552.

\bibitem{PanteleevKalachev}
P.~Panteleev and G.~Kalachev.
\newblock Asymptotically good quantum and locally testable classical LDPC codes.
\newblock In \emph{STOC 2022}, ACM, 2022, 375--388.

\bibitem{SykoraVrto}
O.~S\'ykora and I.~Vrt'o.
\newblock On VLSI layouts of the star graph and related networks.
\newblock \emph{Integration} 17 (1994), 83--93.

\bibitem{Thomassen}
C.~Thomassen.
\newblock Every planar graph is 5-choosable.
\newblock \emph{J. Combin. Theory Ser. B} 62 (1994), 180--181.

\bibitem{Vizing}
V.~G. Vizing.
\newblock Vertex colourings with given colours (in Russian).
\newblock \emph{Diskret. Anal.} 29 (1976), 3--10.

\end{thebibliography}

\end{document}
